\documentclass[11pt,a4paper]{article}
\usepackage[margin=2.3cm]{geometry}
\usepackage{amsmath,amssymb,booktabs,xcolor}
\usepackage[colorlinks=true,linkcolor=blue!50!black]{hyperref}
\title{Finite Element Method: a primer for the Shape Analytics FEM tab}
\author{Shape Analytics project}
\date{}
\begin{document}
\maketitle

\section{Why the tab exists}
The syllabus (Chapters 1--11) derives equilibrium, strain--displacement and stress--strain equations, then solves them in closed form for simple shapes. The 3D tab uses the half-space (Boussinesq/Cerruti) solution, which cannot give zero traction on the other faces of a finite body. The finite element method (FEM) solves the same equations numerically on the real geometry. It is not in the syllabus; the tab is a separate module and can be removed without touching the rest of the app.

\section{The idea in eight steps}
\begin{enumerate}
\item \textbf{Mesh.} Cut the body into small elements joined at nodes (segments, triangles, quadrilaterals).
\item \textbf{Unknowns.} The displacement components at each node, the degrees of freedom $\mathbf{u}$.
\item \textbf{Shape functions.} Inside an element the displacement is interpolated from its nodal values, $u(\mathbf{x})=\sum_i N_i(\mathbf{x})\,u_i$, with $N_i=1$ at node $i$ and $0$ at the others.
\item \textbf{Strain.} Differentiating the interpolation gives the strain, $\boldsymbol{\varepsilon}=\mathbf{B}\,\mathbf{u}_e$ (Chapter~5, small strain).
\item \textbf{Stress.} $\boldsymbol{\sigma}=\mathbf{D}\boldsymbol{\varepsilon}$ with $\mathbf{D}$ the isotropic Hooke matrix in Voigt form (Chapter~6; plane-stress form used in 2D).
\item \textbf{Stiffness.} Minimising the potential energy, or equating internal and external virtual work (Chapter~10), gives the element matrix
\[ \mathbf{k}_e=\int_{V_e}\mathbf{B}^{T}\mathbf{D}\,\mathbf{B}\,dV. \]
It is the force needed at the nodes for a unit nodal displacement.
\item \textbf{Assembly and solution.} Add all $\mathbf{k}_e$ into one global matrix $\mathbf{K}$ by node numbering, apply the supports (fixed displacements) and solve $\mathbf{K}\mathbf{u}=\mathbf{f}$. $\mathbf{K}$ is symmetric (Maxwell--Betti reciprocity, Section~10.3) and singular until rigid-body motion is removed.
\item \textbf{Recovery.} Compute $\boldsymbol\varepsilon$ and $\boldsymbol\sigma$ per element and average at the nodes. Traction on free faces is zero automatically (the \emph{natural} boundary condition), which is the advantage over the half-space formulas.
\end{enumerate}

\section{Elements implemented}
\begin{center}\small
\begin{tabular}{@{}lll@{}}\toprule
Element & Used in & Notes\\\midrule
2-node bar & tapered bar & linear $u$, constant stress; $k=EA_e/h\,[[1,-1],[-1,1]]$\\
Hermite beam & cantilever beam & cubic $v$, unknowns $v,\theta$ per node; exact nodal deflection\\
CST triangle & cantilever, hole & linear $u$, constant strain; stiff in bending\\
Q4 quadrilateral & cantilever, hole & bilinear $u$, $2\times2$ Gauss points, corner extrapolation\\
Hex8 brick & 3D block & trilinear $u$, $2\times2\times2$ Gauss points; solved by conjugate gradients\\\bottomrule
\end{tabular}
\end{center}

\section{Comparison with the syllabus}
\begin{center}\small
\begin{tabular}{@{}lp{9cm}@{}}\toprule
Chapter & Role in FEM\\\midrule
1 & matrix and tensor algebra, index notation\\
2--3 & stress matrix output; equilibrium is the equation being solved; traction boundary condition becomes nodal loads\\
4 & principal stresses, $\tau_{\max}$ computed from the recovered stress\\
5 & the $\mathbf{B}$ matrix is the strain--displacement relation\\
6 & the $\mathbf{D}$ matrix is Hooke's law\\
8--9 & beam theory is the reference for the beam and cantilever cases\\
10 & strain energy and reciprocity give the stiffness matrix and its symmetry\\
11 & von Mises and Tresca applied to the FEM stress field\\\bottomrule
\end{tabular}
\end{center}

\section{Verification}
The solver was checked against closed forms during development (the scripts are not kept in the project).
\begin{center}\small
\begin{tabular}{@{}lp{7.5cm}@{}}\toprule
Test & Result\\\midrule
Patch test, CST and Q4, irregular mesh & uniform stress reproduced to $10^{-8}$\\
Uniform bar & exact\\
Tapered bar, tip displacement & converges as $h^2$ (0.935 at 1 element, 0.9999 at 32)\\
Beam, tip and UDL load & nodal deflection exact to $10^{-9}$\\
Plane-stress cantilever, Q4 & within 1\% of the Timoshenko value on a 12-element depth mesh\\
Coarse CST cantilever & 18\% too stiff, Q4 on the same mesh 4\%\\
Plate with hole, Q4 / CST ($n=16$) & $K_t=3.12/3.08$ against Kirsch 3\\
Block, uniform pressure on the whole face & exact: $\sigma_z=-q$ everywhere, $\sigma_x=\sigma_y=0$, crush $-qH/E$\\
Block, patch load, $L=4$, $H=3$ m & centre stress 5\% above the half-space value (rigid base); side-face $\sigma_x$ about 0.4\% of $q$ against 2.6\% for the half-space solution\\
Equilibrium, all cases & reactions equal applied loads to rounding\\\bottomrule
\end{tabular}
\end{center}
The plate-with-hole $K_t$ lies a few percent above 3 because Kirsch's value is for an infinite plate; the plate in the tab has a width of ten hole radii.

\section{Limits}
Plane stress only (no plane strain or axisymmetric elements yet). Dense linear solver, so meshes are capped near 3600 degrees of freedom. The 3D block is cuboid only (uniform bricks, square patch aligned with the mesh, quarter symmetry) and uses 8-node bricks, which are stiff in bending and for nearly incompressible materials. Surface stresses are extrapolated and the least accurate. The 3D stress-field tab itself still uses the half-space solution.
\end{document}
